
Paragraph 5 theme: Challenges of thresholding in the epoch-structured setting.
In many fatigue-driving studies, EEG recordings are segmented into epochs that correspond to specific task events or analysis windows. This epoch-structured setting introduces additional challenges for threshold-based blink detection. When thresholds are estimated from the entire recording, they may be biased by the presence of clean epochs that contain few or no blinks, leading to an elevated threshold that suppresses blink detection in epochs where blinks are more prevalent. This issue is particularly problematic when some epochs have been pre-rejected due to artifacts or signal loss, as the remaining valid epochs may not be representative of the overall blink activity. As a result, naive application of thresholding in an epoch-structured context can lead to systematic under-detection of blinks, which undermines the reliability of fatigue assessments based on blink metrics. Addressing this challenge requires developing blink detection methods that explicitly account for the epoch structure and can adapt threshold estimation to the specific characteristics of each epoch or subset of epochs, rather than relying on a global threshold derived from the entire recording.
